\section{讨论：证据强度与待解决问题}

直接证据最充分的是 BaZn$_2$Si$_2$O$_7$ 的相变、Ba\slash Sr 固溶和 Zn 位取代，HT\nobreakdash-XRD、单晶 XRD 与热膨胀测量的结果相互印证 \cite{lin1999phase,thieme2015ba1,kerstan2012thermal,thieme2016new}。玻璃析晶方面，成核剂和微结构效应已有系统实验，但各研究的玻璃组成、升温速率和晶粒尺寸并不统一，因此“相稳定性”与“样品宏观热膨胀”不能等同 \cite{thieme2016thermal,vladislavova2018crystallisation}。

第一项关键缺口是 BaO--ZnO--SiO$_2$ 三元相图和热力学函数。目前只能结合二元相平衡、实验室 XRD 和热分析推断相区；理论研究提示必须纳入振动自由能和温度效应 \cite{erlebach2021phase}。第二项缺口是 Ba$_5$Y$_{12}$Zn[O(SiO$_4$)]$_8$ 的直接结构证据。仅凭 Ba、Y、Zn 与硅酸根的共存不能建立同构关系，还需可复核的衍射数据和配位精修。

建议的安全、可证伪实验路径为：先在小规模、封闭式高温炉中采用经审查的坩埚材料制备组成梯度样品；同步采集室温/高温粉末 XRD、TG\nobreakdash-DSC、Raman 光谱和显微成分分布图；对呈现单一衍射相的样品再进行同步辐射或单晶结构解析。上述路径沿用 BaZn$_2$Si$_2$O$_7$ 研究中 HT\nobreakdash-XRD 与热膨胀测量互补的思路 \cite{kerstan2012thermal,thieme2016very}，未验证的富 Y 配方不作为可直接执行的实验方案。

结构描述符有助于缩小候选范围。负热膨胀研究表明，四面体转动、链柔性和配位多面体空隙可作为筛选特征；但描述符只能给出候选，不能替代相平衡实验 \cite{kireeva2022pauling,erlebach2021phase}。
